\pagebreak
\section{Pembagian Tugas Tiap Anggota Kelompok}
Pembagian berikut disusun dari rencana kerja pada \path{SPEC.md} dan riwayat \textit{commit} repositori.
\begin{table}[ht]
    \centering
    \small
    \renewcommand{\arraystretch}{1.4}
    \begin{tabularx}{\textwidth}{@{}p{4.6cm}lX@{}}
        \toprule
        \textbf{Nama Lengkap} & \textbf{NIM} & \textbf{Tugas/Kontribusi} \\
        \midrule
        Samuelson Dharmawan Tanuraharja & 13524001 & SVM \textit{from scratch} dengan SMO, kernel linear, polinomial, dan Gaussian; strategi subsampel dan Platt scaling; perbandingan dengan SVC; visualisasi bidang pemisah. \\
        Edward David Rumahorbo & 13524036 & \textit{Naive Bayes} \textit{from scratch}; \textit{data understanding} dan perbandingan distribusi latih-uji; pembanding \texttt{CategoricalNB}; graf struktur \textit{Naive Bayes}; \textit{cost-sensitive classification}; eksperimen \textit{weighted} KNN dan pohon terpangkas; koordinasi laporan. \\
        Narendra Dharma Wistara M. & 13524044 & \textit{Logistic regression} \textit{from scratch} dengan regularisasi L1, L2, dan Elastic Net; \textit{optimizer} Adam; video kontur \textit{log-loss}; struktur repositori dan antarmuka \texttt{save}/\texttt{load}. \\
        Reinhard Alfonzo Hutabarat & 13524056 & DTL C4.5 \textit{from scratch} dengan atribut numerik, kategorikal, dan nilai hilang; visualisasi pohon dengan parameter \texttt{top\_n}; perapian konfigurasi dan evaluasi notebook. \\
        Miguel Rangga Deardo Sinaga & 13524069 & \textit{Pipeline preprocessing}; modul evaluasi dan \textit{stratified k-fold}; KNN \textit{from scratch}; integrasi notebook; \textit{submission} Kaggle dan pemeriksaan reprodusibilitas. \\
        \bottomrule
    \end{tabularx}
    \caption{Pembagian tugas anggota kelompok G02.}
\end{table}

\renewcommand{\refname}{Daftar Pustaka}
\addcontentsline{toc}{section}{Daftar Pustaka}
\begin{thebibliography}{99}
\bibitem{spesifikasi}
Tim Asisten Laboratorium Inteligensi Artifisial, ``Tugas Besar 2: Supervised Learning, IF3170 Inteligensi Artifisial 2026/2027,'' ver. 1.3, Institut Teknologi Bandung, Bandung, Indonesia, Nov. 2026.

\bibitem{slide-learning}
N. U. Maulidevi, ``Introduction to Learning and Supervised Learning,'' IF3170 Inteligensi Artifisial, Edunex ITB, 2026.

\bibitem{slide-crispdm}
F. Z. Ruskanda, ``CRISP-DM dan Data Understanding,'' IF3170 Inteligensi Artifisial, Edunex ITB, 2026.

\bibitem{slide-dataprep}
E. Ferdian, ``Data Preparation,'' IF3170 Inteligensi Artifisial, Edunex ITB, 2026.

\bibitem{slide-knn}
E. Ferdian, ``KNN,'' IF3170 Inteligensi Artifisial, Edunex ITB, 2026.

\bibitem{slide-bayesnet}
N. U. Maulidevi, ``Bayes Net,'' IF3170 Inteligensi Artifisial, Edunex ITB, 2024.

\bibitem{dtl-id3}
N. U. Maulidevi, ``IF3270 DTL-03: Basic DTL Algorithm (ID3),'' YouTube, 2021. [Daring]. Tersedia: \url{https://www.youtube.com/watch?v=uSTgx9NiEHQ}

\bibitem{dtl-overfit}
N. U. Maulidevi, ``IF3270 Issues in DTL-01: Overfitting,'' YouTube, 2021. [Daring]. Tersedia: \url{https://www.youtube.com/watch?v=hkxC117tEu4}

\bibitem{dtl-continuous}
N. U. Maulidevi, ``IF3270 Issues in DTL-02: Continuous-valued Attribute,'' YouTube, 2021. [Daring]. Tersedia: \url{https://www.youtube.com/watch?v=7x-62Kqj04k}

\bibitem{dtl-missing}
N. U. Maulidevi, ``IF3270 Issues in DTL-03: Missing Attribute Values,'' YouTube, 2021. [Daring]. Tersedia: \url{https://www.youtube.com/watch?v=ZR7l0FYZLF8}

\bibitem{dtl-measures}
N. U. Maulidevi, ``IF3270 Issues in DTL-04: Alternative Measures for Selecting Attribute,'' YouTube, 2021. [Daring]. Tersedia: \url{https://www.youtube.com/watch?v=4_Ueiuchxrk}

\bibitem{khodra-lr-what}
M. L. Khodra, ``Logistic Regression: What \& Why,'' YouTube, 2021. [Daring]. Tersedia: \url{https://www.youtube.com/watch?v=QAKZ_pymEzc}

\bibitem{khodra-lr-sga}
M. L. Khodra, ``Logistic Regression: Stochastic Gradient Ascent,'' YouTube, 2021. [Daring]. Tersedia: \url{https://www.youtube.com/watch?v=O5NSNMQO1yU}

\bibitem{svm-what}
F. Z. Ruskanda, ``SVM: What \& Why?,'' YouTube, 2021. [Daring]. Tersedia: \url{https://www.youtube.com/watch?v=-aNlqGRdQy4}

\bibitem{svm-linear}
F. Z. Ruskanda, ``SVM for Linearly Separable Data,'' YouTube, 2021. [Daring]. Tersedia: \url{https://www.youtube.com/watch?v=RTcwgAlUwlo}

\bibitem{svm-nonlinear}
F. Z. Ruskanda, ``SVM for Non-linearly Separable Data,'' YouTube, 2021. [Daring]. Tersedia: \url{https://www.youtube.com/watch?v=4kK5UJV-LhE}

\bibitem{cunningham-knn}
P. Cunningham and S. J. Delany, ``k-Nearest Neighbour Classifiers: 2nd Edition (with Python examples),'' arXiv:2004.04523, 2020.

\bibitem{dtl-algorithms}
GeeksforGeeks, ``Decision Tree Algorithms,'' GeeksforGeeks. [Daring]. Tersedia: \url{https://www.geeksforgeeks.org/machine-learning/decision-tree-algorithms/}

\bibitem{russell-norvig}
S. J. Russell and P. Norvig, \emph{Artificial Intelligence: A Modern Approach}, 4th ed. Hoboken, NJ, USA: Pearson, 2021.

\bibitem{crisp-dm}
P. Chapman \emph{et al.}, ``CRISP-DM 1.0: Step-by-step data mining guide,'' SPSS Inc., 2000.

\bibitem{quinlan-c45}
J. R. Quinlan, \emph{C4.5: Programs for Machine Learning}. San Mateo, CA, USA: Morgan Kaufmann, 1993.

\bibitem{quinlan-mdl}
J. R. Quinlan, ``Improved use of continuous attributes in C4.5,'' \emph{Journal of Artificial Intelligence Research}, vol. 4, pp. 77--90, 1996.

\bibitem{breiman-cart}
L. Breiman, J. H. Friedman, R. A. Olshen, and C. J. Stone, \emph{Classification and Regression Trees}. Belmont, CA, USA: Wadsworth, 1984.

\bibitem{zou-hastie}
H. Zou and T. Hastie, ``Regularization and variable selection via the elastic net,'' \emph{Journal of the Royal Statistical Society: Series B}, vol. 67, no. 2, pp. 301--320, 2005.

\bibitem{beck-teboulle}
A. Beck and M. Teboulle, ``A fast iterative shrinkage-thresholding algorithm for linear inverse problems,'' \emph{SIAM Journal on Imaging Sciences}, vol. 2, no. 1, pp. 183--202, 2009.

\bibitem{cortes-vapnik}
C. Cortes and V. Vapnik, ``Support-vector networks,'' \emph{Machine Learning}, vol. 20, no. 3, pp. 273--297, 1995.

\bibitem{platt-smo}
J. C. Platt, ``Sequential minimal optimization: A fast algorithm for training support vector machines,'' Microsoft Research, Tech. Rep. MSR-TR-98-14, 1998.

\bibitem{keerthi-smo}
S. S. Keerthi, S. K. Shevade, C. Bhattacharyya, and K. R. K. Murthy, ``Improvements to Platt's SMO algorithm for SVM classifier design,'' \emph{Neural Computation}, vol. 13, no. 3, pp. 637--649, 2001.

\bibitem{platt-prob}
J. C. Platt, ``Probabilistic outputs for support vector machines and comparisons to regularized likelihood methods,'' in \emph{Advances in Large Margin Classifiers}. Cambridge, MA, USA: MIT Press, 1999, pp. 61--74.

\bibitem{lin-platt}
H.-T. Lin, C.-J. Lin, and R. C. Weng, ``A note on Platt's probabilistic outputs for support vector machines,'' \emph{Machine Learning}, vol. 68, no. 3, pp. 267--276, 2007.

\bibitem{friedman-bnc}
N. Friedman, D. Geiger, and M. Goldszmidt, ``Bayesian network classifiers,'' \emph{Machine Learning}, vol. 29, pp. 131--163, 1997.

\bibitem{elkan-cost}
C. Elkan, ``The foundations of cost-sensitive learning,'' in \emph{Proc. 17th Int. Joint Conf. Artificial Intelligence (IJCAI)}, 2001, pp. 973--978.

\bibitem{kingma-adam}
D. P. Kingma and J. Ba, ``Adam: A method for stochastic optimization,'' in \emph{Proc. 3rd Int. Conf. Learning Representations (ICLR)}, 2015.

\bibitem{fan-wss}
R.-E. Fan, P.-H. Chen, and C.-J. Lin, ``Working set selection using second order information for training support vector machines,'' \emph{Journal of Machine Learning Research}, vol. 6, pp. 1889--1918, 2005.

\bibitem{fayyad-irani}
U. M. Fayyad and K. B. Irani, ``Multi-interval discretization of continuous-valued attributes for classification learning,'' in \emph{Proc. 13th Int. Joint Conf. Artificial Intelligence (IJCAI)}, 1993, pp. 1022--1028.

\bibitem{bergstra-random}
J. Bergstra and Y. Bengio, ``Random search for hyper-parameter optimization,'' \emph{Journal of Machine Learning Research}, vol. 13, pp. 281--305, 2012.
\end{thebibliography}

\clearpage
\section{Lampiran}
\subsection{Form Penggunaan \textit{Generative} AI}
\begin{spacing}{1.0}
\begin{center}
    {\large\bfseries SURAT PERNYATAAN PENGGUNAAN AI\par}
\end{center}
\vspace{-0.3em}
\noindent Dengan ini, \sout{saya}/kami*:

\begin{tabular}{@{}l@{\hspace{1.2em}}l@{\hspace{0.8em}}p{0.62\textwidth}@{}}
    Nama/NIM & : & \raggedright\arraybackslash \mbox{Samuelson D.T. (13524001)}, \mbox{Edward D.R. (13524036)}, \mbox{Narendra D.W.M. (13524044)}, \mbox{Reinhard A.H. (13524056)}, \mbox{Miguel R.D.S. (13524069)} \\
    Fakultas/Sekolah & : & Teknik Informatika - STEI \\
    Kode Mata Kuliah & : & IF3170 \\
    Nama Mata Kuliah & : & Inteligensi Artifisial \\
    Judul Tugas/Proposal & : & Tugas Besar 2: \textit{Supervised Learning} \\
\end{tabular}

\vspace{0.6em}
\noindent Menyatakan bahwa \sout{saya}/kami* menggunakan/\sout{tidak menggunakan}* AI dalam pengerjaan maupun penyusunan luaran tugas pada mata kuliah yang tertulis di atas.

\vspace{0.4em}
\noindent Jika Ya, maka alat AI/kecerdasan buatan yang digunakan adalah:

\vspace{0.4em}
\begin{table}[H]
    \centering
    \small
    \renewcommand{\arraystretch}{1.15}
    \begin{tabularx}{\textwidth}{|>{\raggedright\arraybackslash}X|c|>{\centering\arraybackslash}p{2.8cm}|}
        \hline
        \multicolumn{1}{|c|}{\textbf{Lingkup Pekerjaan}} & \textbf{Digunakan?} & \textbf{Tingkat Penggunaan} \\
        \hline
        \textbf{Pemeriksaan Ejaan:} menggunakan tools seperti Grammarly, Paperpal, Deepl, Quillbot, Grammarbot, LanguageTool, ProWritingAid, ChatGPT, Google Gemini, atau sejenisnya. & Ya/\sout{Tidak}* & Rendah \\
        \hline
        \textbf{Pembuatan Teks:} menggunakan tools seperti ChatGPT, Gemini, Paperpal, Copilot, GrammarlyGO, WordAI, WriteSonic, Jasper, Jenni AI, atau sejenisnya. & \sout{Ya}/Tidak* & Tidak Ada \\
        \hline
        \textbf{Bantuan Diskusi dalam Penyusunan Konten:} menggunakan tools seperti ChatGPT, Gemini, Copilot, Perplexity, Jenni AI, Zoom-Companion, atau sejenisnya. & Ya/\sout{Tidak}* & Sedang \\
        \hline
        \textbf{Pembuatan Gambar, Video, dan/atau Grafik:} menggunakan tools seperti Craiyon, DALL-E, Midjourney, Stable Diffusion, Microsoft Designer, Gemini, Canva AI, atau sejenisnya. & \sout{Ya}/Tidak* & Tidak Ada \\
        \hline
        \textbf{Lainnya, Jelaskan:}\vspace{1.6em} & & \\
        \hline
    \end{tabularx}
\end{table}

\noindent \sout{Saya}/Kami* juga menyatakan bahwa setiap penggunaan AI/kecerdasan buatan yang dilakukan pada mata kuliah tersebut telah dinyatakan di dalam surat ini.

\vspace{0.5em}
\hfill\begin{minipage}{6.5cm}
    \raggedleft
    Bandung, 24 November 2026\par
    a.n. Kelompok G02 Nutellabear\par
    \vspace{0.2em}
    \centerline{\includegraphics[height=2.1cm]{ttd_narendra.jpg}}
    \vspace{0.2em}
    Narendra Dharma Wistara M.\par
\end{minipage}
\end{spacing}
